% !TeX root = ../main.tex

% 中文摘要
\begin{abstract}
  
  高能核物理碰撞中的双喷注事件是一个被广泛研究的主题，特别是在量子色动力学（Quantum Chromodynamics，QCD）和夸克-胶子等离子体（Quark–Gluon Plasma，QGP）的研究中。喷注是由高能量粒子（如夸克或胶子）衰变产生的粒子束，通常在粒子物理实验中被观测到。
  
  上世纪60年代，在SLAC上进行的带点轻子-核子深度非弹性散射实验让我们第一次认识到核子是由类点物质--部分子所构成。引入部分子分布函数是为了描述核子内部部分子的动力学变化情况。

  微扰量子色动力学（perturbative Quantum Chromodynamics，简称pQCD）是一种用于研究强相互作用的理论框架，基于量子色动力学（QCD）的微扰展开方法。
  
  


  %（华南）
  在重离子碰撞中，初态冷核物质效应经常和末态热核物质效应混合在一起，并反应在观测量中。精确理解初态效应，特别是原子核部分子分布函数(nPDF) 相对于自由核子分布函数(PDF)的变化，是研究喷注淬火效应的基准。
  由于涉及非微扰动力学，PDF及其核修正通常需要从实验数据中间接提取。同时，由于分布函数及其核修正依赖于部分子的味道、动量分数x、标度Q2，因此选择恰当的观测量对准确提取这些依赖性非常重要。以往的实验数据主要敏感于quake分布函数，对胶子分布的核修正限制较少，LHC的p-Pb碰撞为研究胶子分布提供了新的机遇。
  我们的工作，主要关注p-Pb碰撞中dijet产生，作为PDF核修正特别是胶子分布的探针。为了在观测量层面上更好的分离初始部分子动量分数x和标度Q2，我们着重研究并构造了dijet产生的三重微分截面及其核修正因子$R_pA$，并发现它可以很好的反应出quark和胶子核修正对x和$Q^2$的依赖性。我们还进一步在观测量层面研究了PDF核修正因子随标度的变化。

  %（PRD）
  通过质子和原子核碰撞里双喷注的三微分散射截面成像出部分子分布的核修正在LHC的p-A碰撞里，双喷注的产生提供了原子核里基本部分子分布，特别是胶子分布的宝贵信息。双喷注的三重微分截面使核子中部分子分布函数（nPDF）的运动学扫描得到了很好的控制。
  我们研究了LHC实验上上质子-质子碰撞和质子-铅核（$p$Pb）碰撞中双喷注产生的三重微分截面的几种类型，并在微扰QCD框架下达到了次领头阶。我们计算p-Pb碰撞时用了四个nPDF参数化形式：EPPS16，nCTEQ15，TUJU19，和nIMParton16。
  我们表明了，这类可观测的散射截面的核修正因子 $R_{pPb}$ 可以很好的成像出部分子分布的核修正，后者定量为比值…。我们可以观测和直观地理解这四个nPDF参数化形式所预言的RpPb有很大差异。预计未来对这些可观测量的测量不仅可以限制nPDF的参数化形式，而且还有助于确定各种核效应，如在不同x区域的遮蔽、反遮蔽、EMC和费米运动，以及随探测标度$Q^2$产生的相关变化。

  %
  在相对论性重离子碰撞中会产生一种解离的夸克和胶子等离子体，称为夸克胶子等离子体（QGP）。QGP形成的最显著实验特征之一是在高动量转移的初始非弹性散射中产生的高横动量（$p_T$）喷注的抑制。测量结果表明，在重离子碰撞中，与质子-质子碰撞相比，高（$p_T$）喷注的产生受到强烈抑制。
  对高温高密的核物质性质的研究是核物理的重要课题，而核修正因子是实验室研究核物质性质极其重要的探针。

  %
  在相对论重离子对撞机（Relativistic Heavy-Ion Collider）和大型强子对撞机（Large Hadron Collider）中的高能重离子碰撞已经从定性的理解发展为对量子色动力学介质在极高温度下性质的精确提取。


  %(ATLAS: Xe-Xe)
  对双喷注的压低和关联的测量是通过LHC的ATLAS探测器中的Xe-Xe数据进行收集的。
  使用$R=0.4$ anti-$k_T$算法重建的双喷注在喷注$p_T$的范围从32 GeV到398 GeV以及对撞的中心度上有不同的测量。在最中心的Xe+Xe碰撞中发现了明显的双喷注动量不平衡，而在更边缘的碰撞中则有所减少。

  %（Xe-Xe V2.pdf）
  我们首次对Xe+Xe碰撞中双喷注的动量平衡进行了介质修正的理论研究。在这项研究中，我们采用了POWHEG框架下的次领阶（NLO）QCD双喷注过程的矩阵元，并结合PYTHIA8模拟部分子喷注（PS），以在pp参考中准确描述动量平衡。然后，我们使用一个结合辐射能损和碰撞能损的框架来模拟高能量部子在炽热而密集的介质中的传播过程。理论计算得出的动量平衡结果对于ATLAS在Xe+Xe碰撞中的测量提供了良好的描述。此外，我们还计算了从全包含喷注中选择的c$\bar{c}$和b$\bar{b}$双喷注的动量平衡，并对这三种类型的双喷注进行比较，以研究能量损失的质量和味道依赖性。我们观察到相对于p+p碰撞和Xe+Xe碰撞中非遍举双喷注和来说，平衡c$\bar{c}$和b$\bar{b}$双喷注的比例更高。相对于b$\bar{b}$双喷注，c$\bar{c}$双喷注有几乎同样的动量平衡度。由于在p+p碰撞中，c$\bar{c}$和b$\bar{b}$双喷注的初始动量平衡谱几乎相同，我们可以得出结论，在给定的$p_T$区域内，对于粲夸克和底夸克来说，介质修正效应几乎相同，在Xe+Xe碰撞中高$p_T$区域的能量损失没有明显的质量依赖性。





% 关键词
% 冷核物质效应，部分子分布函数，双喷注，相对论中重离子碰撞，夸克胶子等离子体，喷注淬火，核修正因子





\end{abstract}



% 英文摘要
\begin{abstract*}

\iffalse
An abstract of a dissertation is a summary and extraction of research work and contributions. Included in an abstract should be description of research topic and research objective, brief introduction to methodology and research process, and summary of conclusion and contributions of the research. An abstract should be characterized by independence and clarity and carry identical information with the dissertation. It should be such that the general idea and major contributions of the dissertation are conveyed without reading the dissertation.

An abstract should be concise and to the point. It is a misunderstanding to make an abstract an outline of the dissertation and words ``the first chapter'', ``the second chapter'' and the like should be avoided in the abstract.

Keywords are terms used in a dissertation for indexing, reflecting core information of the dissertation. An abstract may contain a maximum of 5 keywords, with semi-colons used in between to separate one another.

\fi

%  （PRD）
Dijet production in proton-nucleus (pA) collisions at the LHC provides invaluable information on the
underlying parton distributions in nuclei, especially the gluon distributions. Triple-differential dijet cross
sections enable a well-controlled kinematic scan (over momentum fraction x and probing scale $Q^2$) of the
nuclear parton distribution functions (nPDFs),i.e.
In this work,we study several types of triple-differential cross sections for dijet production in proton-proton and proton-lead (pPb) collisions at the
LHC, to next-to-leading order within the framework of perturbative quantum chromodynamics (pQCD).
Four sets of nPDF parametrizations, EPPS16, nCTEQ15, TUJU19, and nIMParton16 are employed in the
calculations for pPb collisions. 
We show that the observable nuclear modification factor RpPb of such cross
sections can serve as a nice image of the nuclear modifications on parton distributions, quantified by the
ratio 
Considerable differences among the R predicted by the four nPDF sets can be observed and intuitively understood. Future measurements of such observables are expected to not only constrain the nPDF parametrizations, but also help confirm various nuclear effects, e.g., shadowing, antishadowing, EMC, and Fermi motion in different regions of x and their variations with probing scale $Q^2$.
$r^{\textrm{A}}_i(x,Q^2)\!=\!f^\textrm{A,proton}_i(x,Q^2)/f^\textrm{proton}_i(x,Q^2)$

%(ATLAS: Xe-Xe)
Measurements of the suppression and correlations of dijets is performed using 3 $\mu b^{-1}$ of Xe+Xe data collected with the ATLAS detector at the LHC.
Dijets with jets reconstructed using the R = 0.4 anti-$k_T$ algorithm are measured differentially in jet $p_T$ over the range of 32 GeV to 398 GeV and the centrality of the collisions. Significant dijets momentum imbalance is found in the most central Xe+Xe collisions, which decreases in more peripheral collisions...
The differences between the dijet suppression in Xe+Xe and Pb+Pb are further quantified by the ratio of pair nuclear-modification factors. The results are found to be consistent with those measured in Pb+Pb data when compared in classes of the same event activity and when taking into account the difference between the center-of-mass energies of the initial parton scattering process in Xe+Xe and Pb+Pb collisions.
These results should provide input for a better understanding of the role of energy density, system size, path length, and fluctuations in the parton energy loss.

\end{abstract*}